% Short running form: the full title is wider than the box the running head
% leaves beside the Sunday's name.
\runninghead{True Words from a Divided Heart}
\properlane{sincere-heart}{True Words from a Divided Heart: The Hypocrites' Question and the One Who Reads Hearts}
\label{sec:sincere-heart}

Read from its middle, the Gospel is a scene about the one who renders before it
is a ruling about what is rendered. Men come to Jesus with a question and
open it with praise, and every word of the praise is true. He answers by
calling them hypocrites. This interpretation asks why, and follows the
witnesses who locate the fault not in the words but in the heart behind them.
Jerome defines the hypocrite as one who is one thing and pretends another; John Chrysostom says that they spoke the truth
``not with an upright mind''; Aquinas says that they had one thing in the
mouth and another in the heart; Augustine says that they spoke one way and
judged another. Chrysostom, Hilary of Poitiers and Aquinas add that Christ's
reading of their heart showed who he is; Jerome states the same discernment
as the first virtue of anyone who answers. Carried through the whole Mass,
that contrast shows the Mass asking for the opposite of the questioners'
duplicity, a will devoted and a heart sincere, and setting the Church in a
crooked generation as lights whose words and lives agree.

\subsection{Why true words are called hypocrisy}

The difficulty is in the Gospel's own sequence. The questioners say, ``Master,
we know that thou art a true speaker and teachest the way of God in truth.
Neither carest thou for any man: for thou dost not regard the person of men''
(Mt 22:16). Every clause is accurate; Matthew's readers know that Jesus is
truthful and teaches God's way, and the rest of the scene shows that he
regards no one's person. Yet two verses later Jesus calls them hypocrites
(22:18). If the charge does not concern the truth of what they said, what
does it concern?

Jerome answers by defining the word from the scene. \latin{Hypocrita ergo
appellatur, qui aliud est, et aliud simulat, id est, aliud opere agit, et
aliud voce praetendit}: ``he is called a hypocrite, then, who is one thing and
feigns another, that is, does one thing in deed and puts forward another with
his voice'' (\work{Commentarii in Matthaeum} III, PL~26, col.~163). On the
praise itself he says, \latin{Blanda et fraudulenta interrogatio}, ``a
flattering and deceitful question'', which draws the one who answers toward
fearing God rather than Caesar, so that he may say tribute should not be paid
(cols.~162--163). The words are true; their purpose is a snare.

John Chrysostom makes the same distinction with more patience. ``And see their
flattery, and their hidden craft.'' They call him truthful, but ``How said ye
then, `He is a deceiver' ''? They ``testify the truth unto Him, confessing what
was really so, nevertheless, not with an upright mind, nor willingly''
(\work{Homilies on Matthew} 70.1). Truth spoken without an upright mind is
still truth; but it is not testimony, because the speaker does not mean it.
Chrysostom also notices the turn of their question: ``They say not, Tell us
what is good, what is expedient, what is lawful? but, `What thinkest Thou?' ''
They do not want to learn; they want something to use.

Aquinas gives the anatomy of the act. \latin{Homines mali ab adulatione
incipiunt}, ``evil men begin from flattery''. He explains why the praise is
nonetheless true of Christ: \latin{Verax est qui veritatem loquitur; et hoc
est proprium Dei, et eius qui Deo coniunctus est \ldots\ Christus autem est
Deo coniunctus per unionem, et ideo verax est}, ``truthful is he who speaks the
truth; and this is proper to God and to one joined to God \ldots\ but Christ
is joined to God by union, and therefore he is truthful''. The questioners'
lie lay elsewhere: they spoke against their own heart, since they held him not
a master but a deceiver. And so: \latin{hypocritae sunt proprie qui aliud in
ore habent, aliud in corde}, ``hypocrites, properly speaking, are those who
have one thing in the mouth, another in the heart'' (\work{Super Matthaeum}
c.~22, Venice ed., p.~283). Augustine, expounding a psalm that bids men speak
justice and judge rightly, puts the same charge in the questioners' own terms:
``Of paying tribute do ye think? To do justice are ye willing? \ldots\ But now
because in one way ye speak, in another way judge, hypocrites ye are''
(\work{Enarrationes in Psalmos} 57.8).

The four witnesses agree on the diagnosis by distinct reasoning. Jerome sets
deed against voice, Chrysostom mind against testimony, Aquinas heart against
mouth, Augustine judgment against speech. None of them makes the fault a matter
of the words' content; all of them place it in the distance between the words
and the one who speaks them. The vice they name belongs to no people and no
party. The Gospel names those who came, disciples of the Pharisees sent with
the Herodians, and the vice can belong to anyone who praises in order to
use.

\subsection{He answers the heart before the words}

Matthew says that Jesus answered ``knowing their wickedness'' (22:18;
\latin{Cognita autem Iesus nequitia eorum}). Under that phrase the witnesses
ask what kind of knowing it was. Hilary of Poitiers says that Christ,
\latin{interna cogitationum secreta cognoscens (Deus enim nihil non eorum quae
intra homines sunt absconsa speculatur)}, ``knowing the inner secrets of their
thoughts, for God sees everything hidden within men'', asked for the denarius
(\work{Commentarius in Matthaeum} 23.1, PL~9, col.~1044). Aquinas divides the
answer in two: \latin{Et primo respondet ad mentem eorum; secundo ad verba
\ldots\ Quia hominis est respondere ad verba, Dei autem respondere ad mentem;
ideo quia Christus erat Deus, et homo, ideo ad utrumque respondet}, ``and first
he answers their mind, then their words \ldots\ because it belongs to man to
answer words, but to God to answer the mind; and therefore, because Christ was
God and man, he answers both'' (p.~283). The rebuke ``ye hypocrites'' is the
answer to their mind; the coin is the answer to their words.

Here the interpretation meets its difficulty. Is the reading of their heart a
sign of who Jesus is, or only the discernment that any wise person should have
when flattered? The objection can draw on Jerome, who states the same
discernment as a rule for every answerer: \latin{Prima virtus est
respondentis, interrogantium mentem cognoscere, et non discipulos, sed
tentatores vocare}, ``the first virtue of one who answers is to know the mind
of those who ask, and to call them not disciples but tempters'' (col.~163).
Jerome draws no conclusion about Christ from this; he gives a rule. The
objection might add that the snare was laid before a crowd, and that a shrewd
hearer could have guessed it. Matthew, however, says only that the Pharisees
took counsel how to ensnare him (22:15) and that Jesus knew their wickedness
(22:18); he does not say whether anyone else knew their purpose beforehand.

Chrysostom's own reasoning answers the objection from the scene itself. The
questioners' words gave nothing away: ``And yet their words were full of much
respect, for they both called Him Master, and bore witness to His truth, and
that He was no respecter of persons; but being God, He was deceived by none of
these things. Wherefore they also ought to have conjectured, that the rebuke
was not the result of conjecture, but a sign of His knowing their secret
thoughts'' (70.1). Because the intent was hidden behind true and respectful
words, the rebuke could not have been a guess. And when the scene is over he
draws the conclusion: ``He gave them proof of His Godhead, by revealing the
secrets of their hearts, and with gentleness did He silence them'' (70.2).

The christological claim is Chrysostom's, Hilary's and Aquinas's, and
Chrysostom grounds it in the concealment. Jerome's sentence is a rule of
prudence for any answerer, a sign of no one's Godhead, and it belongs with the
moral lesson Aquinas draws from the same verse: \latin{datur nobis exemplum,
quod non debemus adulatoribus credere}, ``an example is given us, that we ought not to believe flatterers''
(p.~283).

\subsection{The counter-question: the verdict from their own mouths}

Jesus does not answer the question as asked. He asks for the coin and asks
whose image it bears. Augustine, Chrysostom, Jerome and Aquinas agree that the
counter-question is not an evasion but the means of the answer, and that the
answer comes out of the questioners' own mouths. They give different reasons.

Augustine explains it as a rule of disputation. Sometimes, he says, a
questioner should be asked something else, ``by telling which he would answer
himself''; those who asked about tribute ``were asked another question, viz.
whose image the money bore which was brought forward by themselves; and
because they told what they had been asked \ldots\ they gave a kind of answer
to themselves in reference to the question they had asked the Lord: and
accordingly from their answer He drew this inference, `Render therefore unto
Caesar \ldots' '' He sets this case apart from the question about Jesus'
authority, where the questioners refused to answer and Jesus said, ``Neither
tell I you'' (\work{De sermone Domini in monte} II.20.70). In his exposition
of Psalm~57 he hears the psalm's teeth broken in their own mouth: Jesus ``saith
not, it is Caesar's: but asketh Whose it is? in order that their teeth in their
own mouth might be utterly broken \ldots\ Now ye have made answer''
(\work{Enarrationes in Psalmos} 57.8).

Chrysostom: ``by their tongue He brings out the decision, and causes them to
decide, that it is lawful; which was a clear and plain victory. So that, when
He asks, not from ignorance doth He ask, but because it is His will to cause
them to be bound by their own answers'' (70.2). Jerome: \latin{Sapientia
semper sapienter agit, ut suis potissimum tentatores sermonibus confutentur},
``wisdom always acts wisely, so that tempters are refuted above all by their
own words''; and Jesus asks whose the image is not in ignorance but \latin{ut ad
sermonem eorum competenter respondeat}, ``that he may answer fittingly to
their word'' (col.~163). Aquinas: \latin{non potest melius confutare
opponentem, quam secundum verba sua. Unde primo ponit quaestionem; secundo ex
responsione elicit veritatem}, ``he cannot better refute an opponent than by
his own words; so first he puts the question, then from the answer he draws
out the truth'', asking \latin{ex dispensatione}, by design and not from
ignorance (pp.~283--284).

The reasons differ: the questioner made to answer himself, the teeth broken in
their own mouth, the binding by their own answer, wisdom refuting by the
tempters' words. They converge on one point: the one who reads the heart also makes the hypocrite's own mouth give
the verdict. The divided speaker is answered by the very words he had used to
hide himself.

\subsection{The remedy: a single intention}

If hypocrisy is a divided heart, its contrary is a heart that is one. The
acclamation's verse comes from a sentence that names that contrary. Paul wants
the Philippians to be ``blameless and sincere children of God'', in the Latin
\latin{simplices}, before they shine as lights. Aquinas comments: \latin{Filius
enim est similis patri, Deus autem simplex est: unde simplices simus sicut filii
Dei; quod est, quando intentio est ad unum}, ``for a son is like his father,
and God is simple; so let us be simple as sons of God, which is when the
intention is toward one thing''; and he cites James: ``A double minded man is
inconstant in all his ways'' (Jas 1:8) (\work{Super ad Philippenses} c.~2,
lect.~4). The word ``sincere'' stands in an earlier clause of v.~15, outside the
acclamation's words, which take only the shining and the word of life. But
the light Paul describes is the light of children whose intention is
single.

The second reading shows the same agreement of word and power in the
preaching that founded the Thessalonian church. The gospel came ``not \ldots\
in word only, but in power also: and in the Holy Ghost and in much fulness'',
and Aquinas sets that power against loftiness of speech (\work{Super I ad
Thessalonicenses} c.~1, lect.~1). Chrysostom adds that God permitted
persecution from the first ``that no one might say, that the Preaching was established lightly or by
flattery \ldots\ that it was not human persuasion, but the power of God, that
persuaded the souls of the believers'' (\work{Homilies on First Thessalonians}
1). The questioners in the temple used true words as flattery; the gospel
reached Thessalonica without flattery, and its words and power agreed.

Augustine, expounding the psalm of this Mass, gives a name to the divided
heart: ``What is the whited wall, save hypocrisy, that is, pretence? Without
it is bright, within it is dirt'' (\work{Enarrationes in Psalmos} 95.3).

\subsection{The Missal's texts in this reading}

The Collect asks that the assembly serve God's majesty \latin{sincero corde}.
In this interpretation those words name the precise contrary of what the
Gospel exposes: a heart not divided between what it says and what it intends.
The Collect also asks for a will devoted to God \latin{semper}, always, not
when it is useful.

The Entrance Antiphon is the prayer of one who calls on God because God
hears. Its psalm begins, in verses not sung, with a prayer ``which proceedeth
not from deceitful lips'' and a heart God has proved (Ps 16:1, 3). Augustine,
reading the psalm as the voice of Christ with his Church, glosses the first
verse: ``not going forth to Thee in deceitful lips \ldots\ that I may not in
judging utter aught else than I understand in Thee'' (\work{Enarrationes in
Psalmos} 16.2). Bellarmine, on the verse about the proved heart, says that
``there are many wicked persons, to all appearance with a fair exterior''. The
antiphon's own words are a cry of trust; the sincerity belongs to the psalm
around them.

The first reading contributes only by contrast, and less than any other text.
Cyrus does not know the Lord and yet serves his purpose; Jerome glosses ``thou hast not known me'' as ``you
worshipped images, not God''. The questioners say ``we know'' and yet tempt.
Aquinas says that the fruit God required of Cyrus's power was the knowledge of
God. Set beside the Gospel, the two show that knowing God appears in what is
done and not in what is said. Cyrus is no model of sincerity; he serves a
purpose he does not know.

The psalm bids the nations bring the Lord glory and enter his courts, and
Augustine's present for those courts is a humble heart (95.9). The gift the
psalm asks for is a heart, and it is judged by the Lord who will ``judge the people with justice'' (v.~10c).

The Prayer over the Offerings asks to serve God's gifts \latin{libera mente},
with a free mind, and to be cleansed by the mysteries served. A mind free of
self-interest, purified at the altar, is the opposite of the questioners'
mind, which served its own purpose under a show of respect.

The first Communion antiphon names the eyes of the Lord upon those who fear
him and hope in his mercy: the God who sees. Bellarmine says that fear joined
to hope ``is the mark of real love''. The questioners feigned respect; the
antiphon describes reverence that is real. Where the second antiphon is sung,
it names the Son of Man who gives his life as a ransom. The truthful teacher
whom the questioners praised in order to trap gives his life for many; Bede
reads the verse of one who took the form of a servant.

The Prayer after Communion asks to be instructed, \latin{erudiamur}, by
eternal things: a formation of the heart, which the questioners did not seek.
They came, as Chrysostom observes of their ``What thinkest Thou?'', for an
opinion they could use, not for teaching.

\subsection{Where this interpretation strains}

The interpretation reads every text through a contrast with the questioners,
and that is its difficulty. A contrast drawn against a vice is easily turned
against people, and some of the Fathers who expound this scene pass, on the same
pages, from those who came to Jesus to a whole people. The vice the Gospel
exposes belongs to anyone who praises in order to use, and the Mass's
petitions for a devoted will and a sincere heart are the assembly's prayer
about itself.

\subsection{The four senses of this interpretation}

\begin{fourSenses}
\item[Literal.] The disciples of the Pharisees and the Herodians flatter Jesus
with true words in order to trap him; knowing their wickedness, he calls them
hypocrites, and he answers through the coin, drawing the verdict from their own
mouths (Mt 22:15--21). Hypocrisy is the distance between mouth and heart
(Jerome, Chrysostom, Aquinas, Augustine).
\item[Allegorical.] Christ is the Truth, truthful by his union with God
(Aquinas), who reads hearts as only God does (Hilary, Chrysostom); the Church,
to whom the gospel came in power and not in word only, shines with a borrowed
light in a crooked generation (Aquinas on Philippians).
\item[Moral.] Be single, ``when the intention is toward one thing'' (Aquinas):
a sincere heart and a will devoted always (the Collect); no flattery given or
believed (Aquinas); deeds that match words (Jerome); discernment of those who
praise in order to use (Jerome).
\item[Anagogical.] The Lord comes to judge the peoples with justice (Ps 95:10c)
and, at the psalm's close, the world with justice and the peoples with his
truth (95:13, not sung), when nothing hidden within men will be hidden from the God who, Hilary says, sees
all that lies concealed within them.
\end{fourSenses}
